\documentclass[11pt]{article}
\usepackage{mathblatt}
\usepackage{multicol}


\begin{document}
\pfheftstil
\pfheftkopf{Gleichungssysteme \textperiodcentered{} Lösungen}{}
\begin{pfloesung}{}
\lz{1.}{3x; 2x + 2y}{}{}
\end{pfloesung}
\begin{pfloesung}{}
\lz{2.}{20 $-$ 2y + 3y = 26 $\Rightarrow$ y = 6}{}{}
\end{pfloesung}
\begin{pfloesung}{}
\lz{3.}{von beiden}{}{}
\end{pfloesung}
\begin{pfloesung}{}
\lz{4.}{vor $x$ und $y$: $2$, $5$ und $3$, $1$}{rechts: 58\,€ und 35\,€; in €}{}
\end{pfloesung}
\begin{pfloesung}{}
\lz{5.}{$y$: $5$, $4$, $3$, $2$}{}{}
\end{pfloesung}
\begin{pfloesung}{}
\lz{6.}{$69$}{I: 21, II: 2x + 5y}{}
\end{pfloesung}
\begin{pfloesung}{}
\lz{7.}{$7$}{I: 31, II: x - y}{}
\end{pfloesung}
\begin{pfloesung}{}
\lz{8.}{$28$}{I: 42, II: x = 2y; 14, x}{}
\end{pfloesung}
\begin{pfloesung}{}
\lz{9.}{$118{,}40$}{I: 120, II: 0{,}80x + 1{,}20y}{}
\end{pfloesung}
\begin{pfloesung}{}
\lz{10.}{2x + 2y = 130; x + 3y = 105; Erwachsene 45 €, Kinder 20 €}{}{}
\end{pfloesung}
\begin{pfloesung}{}
\lz{11.}{7 Drei-Bett-Zimmer; 9 Fünf-Bett-Zimmer}{16; 3x + 5y $\Rightarrow$ 66; 3(16 $-$ y) + 5y $\Rightarrow$ 66 = 9}{}
\end{pfloesung}
\begin{pfloesung}{}
\lz{12.}{I: 2x + 2y = 77,80, II: x + 3y = 64,90; Erwachsener 25,90 €, Kind 13,00 €}{64,90 $-$ 3y; 129,80 $-$ 4y $\Rightarrow$ 77,80 = 13}{}
\end{pfloesung}
\begin{pfloesung}{}
\lz{13.}{Gegenzahlen}{}{}
\end{pfloesung}
\begin{pfloesung}{}
\lz{14.}{Q(2; $-$6)}{}{}
\end{pfloesung}
\begin{pfloesung}{}
\lz{15.}{b = $-$4, c = $-$5; y = x² $-$ 4x $-$ 5}{}{}
\end{pfloesung}
\begin{pfloesung}{}
\lz{16.}{x = 16, y = 8}{}{}
\end{pfloesung}
\begin{pfloesung}{}
\lz{17.}{7x = 35 $\Rightarrow$ x = 5}{}{}
\end{pfloesung}
\begin{pfloesung}{}
\lz{18.}{2x = $-$10 $\Rightarrow$ x = $-$5}{}{}
\end{pfloesung}
\begin{pfloesung}{}
\lz{19.}{S(1,5; 1)}{}{}
\end{pfloesung}
\begin{pfloesung}{}
\lz{20.}{II $-$ I: 2x = 3 $\Rightarrow$ x = 1,5; y = 6}{}{}
\end{pfloesung}
\begin{pfloesung}{}
\lz{21.}{6 = $-$9 + 6b + c und 3 = $-$16 + 8b + c $\Rightarrow$ b = 2, c = 3; y = $-$0,25x² + 2x + 3}{}{}
\end{pfloesung}
\begin{pfloesung}{}
\lz{22.}{Apfelbaum 27,20 €, Pflaumenbaum 35,80 €}{63 $-$ x; 6x = 163,20}{}
\end{pfloesung}
\begin{pfloesung}{}
\lz{23.}{$x$: Preis eines Blocks in €, $y$: Preis eines Füllers in €}{}{}
\end{pfloesung}
\begin{pfloesung}{}
\lz{24.}{$x$: Preis eines Zelts in €, $y$: Preis eines Schlafsacks in €}{}{}
\end{pfloesung}
\begin{pfloesung}{}
\lz{25.}{$230$: Einnahmen zusammen in €}{8: 8: Preis eines großen Glases in €}{}
\end{pfloesung}
\begin{pfloesung}{}
\lz{26.}{x: Preis eines Apfelbaums in €, y: Preis eines Pflaumenbaums in €}{Koeffizienten: 9 und 3}{}
\end{pfloesung}
\begin{pfloesung}{}
\lz{27.}{Im Korb liegen zusammen $25$ Äpfel und Birnen}{}{}
\end{pfloesung}
\begin{pfloesung}{}
\lz{28.}{$x$: Preis einer Tischtennisplatte in €, $y$: Preis eines Kickers in €}{}{}
\end{pfloesung}
\begin{pfloesung}{}
\lz{29.}{5n $-$ 7 = 38 $\Rightarrow$ 5n = 45 $\Rightarrow$ n = 9}{}{}
\end{pfloesung}
\begin{pfloesung}{\textbf{Prüfstein}}
\lz{a)}{I. x + y = 3,80; II. 6x + 5y = 21,20}{x Preis Rose, y Preis Tulpe}{}
\lz{b)}{Filip kauft insgesamt 13 Blumen (Rosen und Tulpen zusammen); r = 8 Rosen, t = 5 Tulpen}{13 $-$ t; 29,9 $-$ 0,6t $\Rightarrow$ 26,9 = 0,6t = 3}{}
\end{pfloesung}
\end{document}